\documentclass[10pt]{article}
\usepackage[margin=20mm]{geometry}
\usepackage{amsmath,amssymb,fontspec,unicode-math,xurl,hyperref}
\IfFontExistsTF{DejaVu Serif}{\setmainfont{DejaVu Serif}}{\setmainfont{Latin Modern Roman}}
\setmathfont{latinmodern-math.otf}
\setlength{\parindent}{0pt}\setlength{\parskip}{6pt}\sloppy
\begin{document}
Abstract: The binary probe used in this study cannot establish the asserted association between visible usernames and hidden user IDs: an exact synthetic construction attains perfect accuracy without storing any such association. The supporting mechanism also uses an upper probability bound as a lower bound and assumes a zero expected residual gradient that does not follow from task irrelevance. Independent finite calculations and a locally repaired Gaussian argument delimit these Major findings. The reported representation changes remain separate evidence, and scientific adjudication is pending.\par
Author: AI-assisted technical review\par
Date:\par
Title: Technical Verification: Watch Out Your Album! On the Inadvertent Privacy Memorization in Multi-Modal Large Language Models\par
\# Paper identity and reviewed version\par
This report concerns Watch Out Your Album! On the Inadvertent Privacy Memorization in Multi-Modal Large Language Models, by Tianjie Ju et al., ICML 2025, PMLR volume 267; identifier `I25-43674`, OpenReview `vSzcsYucCW`. The reviewed object is the frozen published proceedings PDF, retrieved on 25 October 2026. Page references below are one-based PDF pages, which coincide with the internal page numbers. They are not the continuous PMLR volume pagination.\par
\# Scope and operational definitions\par
The work inserts synthetic usernames and identifiers into otherwise unrelated VQA images, then studies training-gradient changes and binary probes of hidden representations. This review covers the main method and results, Appendix A’s probability argument, Appendices B–E’s splits/content/transformations/ablation, and static code at commit . Source pages 7, 8, 14 and 15 were visually inspected. No model, author training program, privacy extraction or attack was executed. All independent examples below are artificial finite calculations.\par
The principal distinction is among (i) sensitivity to an input watermark, (ii) information about whether a username belongs to a previously exposed set, and (iii) retention of the exact association between that username and an unobserved identifier. These are different estimands. The paper uses the first two measurements to support claims about the third.\par
\# The probe cannot identify the hidden association\par
Problem 1 (Binary set classification is insufficient to establish username–ID memory).  Section 4.3.1, Scenario II, page 6, concludes that the model internally links usernames to their appropriate user IDs from a high binary probing accuracy. The same inference appears in Section 4.3.2, page 7. The measured labels do not test that association.\par
Sections 3 and 4.1 define two fixed sets \(U_1,U_2\), each with five usernames. The probing image displays the username, and the target is whether it belongs to \(U_1\) or \(U_2\). It is not the identifier value. Splitting images into probe training and test subsets does not separate these ten underlying identities. Identical generation settings for the two sets make them exchangeable before sampling, but do not make the fixed realized usernames indistinguishable to a classifier trained on repeated examples of them.\par
\#\# A constructive identifiability check\par
Number the ten artificial usernames \(j=1,\ldots,10\) and let their hidden representation be \(h(j)=e_j\in\mathbb R^{10}\), obtained solely from the visible username. Let the binary label be \(b(j)=0\) for \(j\leq5\) and \(1\) otherwise. A linear classifier with weights \((-1,-1,-1,-1,-1,1,1,1,1,1)\) has perfect accuracy. Neither \(h\) nor the classifier contains a user ID or any username–ID pairing.\par
Now permute the five artificial ID values paired with the five seen usernames in all \(5!=120\) possible ways. The displayed usernames, representations and binary predictions remain identical under every permutation. A fixed arbitrary guess of the associations has only \(1/5\) mean correct recovery over those permutations, while the binary probe retains 100\% accuracy. The independent script enumerates these assignments. This is a countermodel to the inference, not a claim that the actual MLLM has one-hot representations or zero association memory.\par
The pre-fine-tuning controls in Figures 4–5 already achieve 85.5\% and 82.0\%, respectively. The paper acknowledges that visible usernames explain much of that performance. The subsequent increases to 91.3–92.2\% and 90.0–91.1\% are evidence that the representation or its linear separability changes after exposure. They do not identify which information changed: improved recognition of the same visible names, learned set membership, and exact ID association remain observationally different possibilities that the binary endpoint does not distinguish. A question mentioning an ID does not change the probe’s target into an ID-recovery test.\par
\#\# Code trace and consequence\par
The archived , lines 34–38, derives labels solely from the filename containing `test1` or `test2`. Lines 69–72 and the analogous later blocks use a stratified image-level 80/20 split and logistic regression. No target ID is used. Thus the static implementation supports the measurement interpretation above. This released split/classifier setup differs from the paper’s stated 60/20/20, Adam, ten-epoch linear-probe protocol; equivalence to the historical run is not established.\par
The supported assessment is Major for the claimed hidden-association conclusion, an important part of RQ2. Evidence of changed representations and possible username membership retention survives. The paper’s own zero direct ID-answer accuracy does not prove an absence of latent memory, but it also supplies no independent positive evidence for the association. To establish that association, a controlled evaluation must distinguish correct from permuted username–ID pairings, use genuinely independent synthetic assignments, and compare models with the same username exposure but differing pairing exposure. Such an experiment is proposed, not performed here; the actual latent association remains Unverifiable from the supplied binary measurements.\par
\# The theoretical mechanism is not established by its proof\par
Problem 2 (An upper-tail upper bound is used to infer a lower probability).  Appendix A, pages 14–15, applies Chebyshev to bound the probability of a large sample cross-covariance from above, then asserts that selecting a threshold ensures a probability bounded away from zero. This implication does not follow.\par
For independent Gaussian scalar coordinates with positive variances, \(B>1\) and \(m=B-1\), the centered sample cross-covariance \(S\) has \[\mathbb E S=0,\qquad \mathbb E S^2=\frac{\sigma_y^2\sigma_u^2}{m}.\] For example, write \(S=y^TPu/m\), where \(P=I-\mathbf1\mathbf1^T/B\) is an orthogonal rank-\(m\) projection. Independence gives the variance from \(\operatorname{tr}(P^2)=m\). In Equation (9), the variance of a centered product is written as the uncentered variance product; the correct factor is \((1-1/B)^2\). Its final loose upper bound \(3\sigma_y^2\sigma_u^2/m\) nevertheless remains valid. Equation (10) is an upper bound and cannot prove a nonzero lower bound merely by choosing \(t\) proportional to the standard deviation.\par
The Gaussian phenomenon can be repaired, which is material counterevidence to treating the qualitative idea as false. After standardizing variances, \(S=m^{-1}\sum_{k=1}^m a_kb_k\) with independent standard normal variables. Its fourth moment is \(3(m+2)/m^3\). Applying Cauchy–Schwarz to the part of \(S^2\) above a fraction \(\theta\in(0,1)\) of its mean gives \[\Pr\{|S|>\sqrt{\theta/m}\}\geq
\frac{(1-\theta)^2m}{3(m+2)}.\] Indeed, \(\mathbb E[S^2\mathbf1_{S^2>\theta\mathbb ES^2}]\geq(1-\theta)\mathbb ES^2\), and squaring the Cauchy–Schwarz bound proves the display. This is an actual lower bound under explicit nondegeneracy assumptions. Independent seeded Gaussian calculations at \(B=2,4,8,32\) check the variance scale and give empirical tail frequencies above this analytic bound. The computation illustrates the repaired Gaussian claim, not the distribution of language embeddings. A central-limit appeal alone does not establish that the real vectors satisfy these Gaussian assumptions; the sample covariance is also undefined at \(B=1\).\par
\#\# Independence does not imply zero residual gradient\par
There is a second, more direct gap in Section 2, Equation (3), page 3. Define the residual gradient as the difference between losses with and without a nuisance input. Independence of nuisance \(u\) and target \(y\) does not imply that the residual’s population mean is zero at an arbitrary model parameter.\par
For an exact four-case example, let \(u,y\) be independent uniform signs, let the original input be the constant zero (so \(u\) is independent of both the original input and \(y\)), and use \(f_\theta(u)=\theta u\) with loss \((\theta u-y)^2/2\). At \(\theta=1\), the four equally likely residual gradients are \(0,2,2,0\), so their expectation is one, despite zero population covariance. In general the expectation is \(\theta\). Gradient descent therefore removes the nuisance coefficient toward zero; a nonzero gradient difference can accompany suppression of irrelevant information. It is not by itself a diagnostic of storing its content. This example satisfies task irrelevance and requires neither a small batch nor a spurious empirical correlation.\par
Together these are Major defects in the stated supporting mechanism, not a proof that mini-batches cannot create spurious covariance or that MLLMs cannot memorize. The repaired Gaussian argument proves only covariance fluctuations. Linking those fluctuations to persistent content memory and the particular batch-size observations requires additional model/loss analysis or controlled evidence. Neither the local repair nor the synthetic counterexample establishes an outcome for the reported MLLMs.\par
\# Empirical and reproducibility boundaries\par
The gradient comparisons in Tables 3 and 6 show changes under watermarking and increasing cosine similarity with larger batches. They remain valid descriptive results as reported; the proposed causal interpretation is the issue. The images in Figures 4–6 show separability changes consistent with those descriptions, without resolving the ID-association ambiguity. Direct prompting reportedly has zero ID-answer accuracy. Table 5 actually gives after-minus-before AUC changes of \(-0.8,-0.5,-0.3\) percentage points, although the prose calls them a marginal increase. This small directional error does not alter the reported lack of a strong positive MIA signal.\par
For reproduction, , lines 84–85, indexes `[:, -2, :]` from the last hook output while the paper describes the final token; the intended token position depends on the exact chat prompt and model implementation and was not resolved by running them. These released-protocol discrepancies are recorded as limits rather than attributed automatically to the historical experiments. The data-generation and split utilities were read selectively; all model internals, original hidden-state dumps and modified gradient-saving library code were not audited. The source snapshot excludes datasets and weights by design.\par
\# Conclusion\par
The strong user-ID association claim is not identified by the binary endpoint, and the theoretical mechanism contains both a reversed-bound inference and a false generic mean-gradient premise. These findings limit the claimed explanation without refuting the narrower evidence of input sensitivity and changed representations. Exact synthetic checks and a local Gaussian proof repair are complete. Association-specific model experiments, historical implementation equivalence, author adoption and independent adjudication remain unresolved. Focused independent AI acceptance has checked both Problems and the local subfindings; scientific adjudication remains pending and no human review has occurred.\par
1 Tianjie Ju et al., Watch Out Your Album! On the Inadvertent Privacy Memorization in Multi-Modal Large Language Models, PMLR 267 (2025), published proceedings version. https://proceedings.mlr.press/v267/ju25a.html.\par
\# Evidence and reproducibility\par
The exact source PDF has SHA-256\par
The matching `verification/0305` directory contains the source text extraction, independently written calculation script and actual output, inspected page images, coverage ledger, findings and provenance. Source code supplied by the authors was not executed. The calculation is an independent finite or analytic diagnostic, not a reproduction of the original training experiments. The final PDF is compiled from this TeX; the compiler log and visual receipt identify the exact bytes inspected. This initial review was performed by the delegated AI unit with GPT-6 Astra / Max formally selected. A separate AI unit subsequently checked the stated Problems against the frozen original source, assumptions, decisive arguments, counterevidence and impact boundaries. Its fresh evidence and exact artifact bindings are retained in the same verification directory. No human review has occurred; scientific adjudication is pending and author adoption is unknown.\par
\end{document}
